\chapter*{Use of generative AI}
\phantomsection
\addcontentsline{toc}{plainchapter}{Use of generative AI}

Generative AI tools were used in this thesis for writing assistance, language refinement, and programming support. ChatGPT 5.5 Extended Thinking by OpenAI (\href{https://www.chatgpt.com}{www.chatgpt.com}) was used to improve wording, structure, clarity, and scientific phrasing, particularly in the Introduction, Background and Related Work, Abstract, and Conclusion. It was not used as a single source of information. The research problem, methodology, experiments, analysis, and conclusions were developed by the author, and all AI-generated content was carefully reviewed, verified, and edited before adding it to the thesis document.

Generative AI was also used for programming support. The ChatGPT Codex 5.4 and 5.5 extension in Visual Studio Code assisted with the implementation of the PatchTST training and evaluation framework, as well as analysis and logic of plotting scripts. The robot and dataset recording scripts were implemented mainly by the author. The author gave clear and detailed instructions and requirements, with feedback and improvements after the AI-generated code was reviewed and tested. Decisions regarding analysis, plotting, interpretation of results, and thesis presentation were all made by the author.

Generative AI was only used as a technical and language support tool.